%%%%%%%%%%%%%%%%
%%% lev 23 %%%
%%%%%%%%%%%%%%%%

\Chapter{23}

\Verse{1}
{
	\hP{וַיְדַבֵּר יְהֹוָה אֶל מֹשֶׁה לֵּאמֹר׃}
}
{
	\eP{23 And YHWH spoke to Moses, saying,}
}
{
}

\Verse{2}
{
	\hP{דַּבֵּר אֶל בְּנֵי יִשְׂרָאֵל וְאָמַרְתָּ אֲלֵהֶם מוֹעֲדֵי יְהֹוָה אֲשֶׁר תִּקְרְאוּ אֹתָם מִקְרָאֵי קֹדֶשׁ אֵלֶּה הֵם מוֹעֲדָי׃}
}
{
	\eP{
		“Speak to the children of Israel, and you shall say to
		them: YHWH’s appointed times, which you shall call holy
		assemblies—these are my appointed times:
	}
}
{
%	\footnote{
%		these are my appointed times. The laws of the holidays as expressed in
%		Leviticus reflect a concern with order and arrangement as well. YHWH
%		tells Moses to tell the people what the nation’s holidays are to be,
%		specifically: the Sabbath, the holiday of unleavened bread (maôt), an
%		unnamed holiday related to firstfruits, an unnamed holiday involving
%		horn blowing (known subsequently in Judaism, but not biblically, as the
%		New Year), the Day of Atonement, and the holiday of booths (sukkôt). In
%		each case the text gives information on how to celebrate the holiday. In
%		all but one case it gives the date on which the holiday is to be
%		celebrated. Interestingly, it does not make historical associations for
%		the holidays, not even connecting the feast of unleavened bread to the
%		exodus from Egypt, or the Sabbath to God’s rest at the end of the
%		creation. (The single exception is an instruction in 23:39–43 that the
%		Israelites should live in booths during the feast of sukkôt in order to
%		remember that YHWH had Israel live in booths when He brought them out of
%		Egypt.) It appears that the main concern, as with the food and purity
%		laws, is not to explicate the background and purposes of the
%		holidays—these things are a concern of Genesis, Exodus, Numbers, and
%		Deuteronomy—but rather to instruct the people on how to observe them
%		properly. In the case of the holidays, the concern of Leviticus with
%		categorization and organization is applied to the temporal as well as
%		the spatial, thus institutionalizing the sanctification of time more
%		than in any text thus far. Organization through ritual is thus designed
%		to pervade the people’s life.
%	}
}

\Verse{3}
{
	\hP{שֵׁשֶׁת יָמִים תֵּעָשֶׂה מְלָאכָה וּבַיּוֹם הַשְּׁבִיעִי שַׁבַּת שַׁבָּתוֹן מִקְרָא קֹדֶשׁ כּל מְלָאכָה לֹא תַעֲשׂוּ שַׁבָּת הִוא לַיהֹוָה בְּכֹל מוֹשְׁבֹתֵיכֶם׃}
}
{
	\eP{
		Six days work shall be done, and in the seventh day is a
		Sabbath, a ceasing, a holy assembly. You shall not do any
		work. It is a Sabbath to YHWH in all your homes.
	}
}
{
}

\Verse{4}
{
	\hP{אֵלֶּה מוֹעֲדֵי יְהֹוָה מִקְרָאֵי קֹדֶשׁ אֲשֶׁר תִּקְרְאוּ אֹתָם בְּמוֹעֲדָם׃}
}
{
	\eP{
		“These are YHWH’s appointed times, holy assemblies that
		you shall proclaim at their appointed time:
	}
}
{
}

\Verse{5}
{
	\hP{בַּחֹדֶשׁ הָרִאשׁוֹן בְּאַרְבָּעָה עָשָׂר לַחֹדֶשׁ בֵּין הָעַרְבָּיִם פֶּסַח לַיהֹוָה׃}
}
{
	\eP{
		In the first month, on the fourteenth of the month,
		‘between the two evenings,’ is YHWH’s Passover.
	}
}
{
}

\Verse{6}
{
	\hP{וּבַחֲמִשָּׁה עָשָׂר יוֹם לַחֹדֶשׁ הַזֶּה חַג הַמַּצּוֹת לַיהֹוָה שִׁבְעַת יָמִים מַצּוֹת תֹּאכֵלוּ׃}
}
{
	\eP{
		And on the fifteenth day of this month is YHWH’s Festival
		of Unleavened Bread. You shall eat unleavened bread seven
		days.
	}
}
{
}

\Verse{7}
{
	\hP{בַּיּוֹם הָרִאשׁוֹן מִקְרָא קֹדֶשׁ יִהְיֶה לָכֶם כּל מְלֶאכֶת עֲבֹדָה לֹא תַעֲשׂוּ׃}
}
{
	\eP{
		On the first day you shall have a holy assembly. You shall
		not do any act of work.
	}
}
{
}

\Verse{8}
{
	\hP{וְהִקְרַבְתֶּם אִשֶּׁה לַיהֹוָה שִׁבְעַת יָמִים בַּיּוֹם הַשְּׁבִיעִי מִקְרָא קֹדֶשׁ כּל מְלֶאכֶת עֲבֹדָה לֹא תַעֲשׂוּ׃}
}
{
	\eP{
		And you shall bring forward an offering by fire to YHWH
		for seven days. On the seventh day you shall have a holy
		assembly. You shall not do any act of work.”
	}
}
{
}

\Verse{9}
{
	\hP{וַיְדַבֵּר יְהֹוָה אֶל מֹשֶׁה לֵּאמֹר׃}
}
{
	\eP{And YHWH spoke to Moses, saying,}
}
{
}

\Verse{10}
{
	\hP{דַּבֵּר אֶל בְּנֵי יִשְׂרָאֵל וְאָמַרְתָּ אֲלֵהֶם כִּי תָבֹאוּ אֶל הָאָרֶץ אֲשֶׁר אֲנִי נֹתֵן לָכֶם וּקְצַרְתֶּם אֶת קְצִירָהּ וַהֲבֵאתֶם אֶת עֹמֶר רֵאשִׁית קְצִירְכֶם אֶל הַכֹּהֵן׃}
}
{
	\eP{
		“Speak to the children of Israel, and you shall say to
		them: When you will come to the land that I am giving to
		you and you reap its harvest, you shall bring the first
		sheaf of your harvest to the priest,
	}
}
{
}

\Verse{11}
{
	\hP{וְהֵנִיף אֶת הָעֹמֶר לִפְנֵי יְהֹוָה לִרְצֹנְכֶם מִמּחֳרַת הַשַּׁבָּת יְנִיפֶנּוּ הַכֹּהֵן׃}
}
{
	\eP{
		and he shall elevate the sheaf in front of YHWH for
		acceptance for you. The priest shall elevate it on the day
		after the Sabbath.
	}
}
{
}

\Verse{12}
{
	\hP{וַעֲשִׂיתֶם בְּיוֹם הֲנִיפְכֶם אֶת הָעֹמֶר כֶּבֶשׂ תָּמִים בֶּן שְׁנָתוֹ לְעֹלָה לַיהֹוָה׃}
}
{
	\eP{
		And on the day of your elevating the sheaf, you shall do
		an unblemished lamb in its first year as a burnt offering
		to YHWH,
	}
}
{
}

\Verse{13}
{
	\hP{וּמִנְחָתוֹ שְׁנֵי עֶשְׂרֹנִים סֹלֶת בְּלוּלָה בַשֶּׁמֶן אִשֶּׁה לַיהֹוָה רֵיחַ נִיחֹחַ וְנִסְכֹּה יַיִן רְבִיעִת הַהִין׃}
}
{
	\eP{
		and its grain offering shall be two-tenths of a measure of
		fine flour mixed with oil, an offering by fire to YHWH, a
		pleasant smell, and its libation shall be wine, a fourth
		of a hin.
	}
}
{
}

\Verse{14}
{
	\hP{וְלֶחֶם וְקָלִי וְכַרְמֶל לֹא תֹאכְלוּ עַד עֶצֶם הַיּוֹם הַזֶּה עַד הֲבִיאֲכֶם אֶת קרְבַּן אֱלֹהֵיכֶם חֻקַּת עוֹלָם לְדֹרֹתֵיכֶם בְּכֹל מֹשְׁבֹתֵיכֶם׃}
}
{
	\eP{
		And you shall not eat bread or parched or fresh grain
		until that very day, until you have brought your {\God}’s
		offering. It is an eternal law through your generations in
		all your homes.
	}
}
{
}

\Verse{15}
{
	\hP{וּסְפַרְתֶּם לָכֶם מִמּחֳרַת הַשַּׁבָּת מִיּוֹם הֲבִיאֲכֶם אֶת עֹמֶר הַתְּנוּפָה שֶׁבַע שַׁבָּתוֹת תְּמִימֹת תִּהְיֶינָה׃}
}
{
	\eP{
		And you shall count from the day after the Sabbath, from
		the day of your bringing the sheaf for an elevation
		offering, seven Sabbaths. They shall be complete.
	}
}
{
%	\footnote{
%		seven Sabbaths. Meaning: seven full weeks.
%	}
}

\Verse{16}
{
	\hP{עַד מִמּחֳרַת הַשַּׁבָּת הַשְּׁבִיעִת תִּסְפְּרוּ חֲמִשִּׁים יוֹם וְהִקְרַבְתֶּם מִנְחָה חֲדָשָׁה לַיהֹוָה׃}
}
{
	\eP{
		You shall count until the day after the seventh Sabbath:
		fifty days. And you shall bring forward a new grain
		offering to YHWH.
	}
}
{
}

\Verse{17}
{
	\hP{מִמּוֹשְׁבֹתֵיכֶם תָּבִיאּוּ לֶחֶם תְּנוּפָה שְׁתַּיִם שְׁנֵי עֶשְׂרֹנִים סֹלֶת תִּהְיֶינָה חָמֵץ תֵּאָפֶינָה בִּכּוּרִים לַיהֹוָה׃}
}
{
	\eP{
		From your homes you shall bring bread for an elevation
		offering: two of them. They shall be two-tenths of a
		measure of fine flour, baked with leaven, firstfruits to
		YHWH.
	}
}
{
}

\Verse{18}
{
	\hP{וְהִקְרַבְתֶּם עַל הַלֶּחֶם שִׁבְעַת כְּבָשִׂים תְּמִימִם בְּנֵי שָׁנָה וּפַר בֶּן בָּקָר אֶחָד וְאֵילִם שְׁנָיִם יִהְיוּ עֹלָה לַיהֹוָה וּמִנְחָתָם וְנִסְכֵּיהֶם אִשֵּׁה רֵיחַ נִיחֹחַ לַיהֹוָה׃}
}
{
	\eP{
		And you shall bring forward with the bread seven lambs,
		unblemished one-year-olds, and one bull of the cattle and
		two rams. They shall be a burnt offering to YHWH, with
		their grain offering and their libations, an offering by
		fire, a pleasant smell to YHWH.
	}
}
{
}

\Verse{19}
{
	\hP{וַעֲשִׂיתֶם שְׂעִיר עִזִּים אֶחָד לְחַטָּאת וּשְׁנֵי כְבָשִׂים בְּנֵי שָׁנָה לְזֶבַח שְׁלָמִים׃}
}
{
	\eP{
		And you shall do one goat as a sin offering and two
		one-year-old lambs as a peace-offering sacrifice.
	}
}
{
}

\Verse{20}
{
	\hP{וְהֵנִיף הַכֹּהֵן אֹתָם עַל לֶחֶם הַבִּכֻּרִים תְּנוּפָה לִפְנֵי יְהֹוָה עַל שְׁנֵי כְּבָשִׂים קֹדֶשׁ יִהְיוּ לַיהֹוָה לַכֹּהֵן׃}
}
{
	\eP{
		And the priest shall elevate them with the bread of
		firstfruits, an elevation offering in front of YHWH with
		the two sheep. They shall be a holy thing to YHWH: for the
		priest.
	}
}
{
}

\Verse{21}
{
	\hP{וּקְרָאתֶם בְּעֶצֶם הַיּוֹם הַזֶּה מִקְרָא קֹדֶשׁ יִהְיֶה לָכֶם כּל מְלֶאכֶת עֲבֹדָה לֹא תַעֲשׂוּ חֻקַּת עוֹלָם בְּכל מוֹשְׁבֹתֵיכֶם לְדֹרֹתֵיכֶם׃}
}
{
	\eP{
		And you shall proclaim on that very day: you shall have a
		holy assembly. You shall not do any act of work: an
		eternal law in all your homes through your generations.
	}
}
{
}

\Verse{22}
{
	\hP{וּבְקֻצְרְכֶם אֶת קְצִיר אַרְצְכֶם לֹא תְכַלֶּה פְּאַת שָׂדְךָ בְּקֻצְרֶךָ וְלֶקֶט קְצִירְךָ לֹא תְלַקֵּט לֶעָנִי וְלַגֵּר תַּעֲזֹב אֹתָם אֲנִי יְהֹוָה אֱלֹהֵיכֶם׃}
}
{
	\eP{
		And when you reap your land’s harvest, you shall not
		finish your field’s corner when you harvest, and you shall
		not gather your harvest’s gleaning. You shall leave them
		for the poor and for the alien. I am YHWH, your {\God}.”
	}
}
{
}

\Verse{23}
{
	\hP{וַיְדַבֵּר יְהֹוָה אֶל מֹשֶׁה לֵּאמֹר׃}
}
{
	\eP{And YHWH spoke to Moses, saying,}
}
{
}

\Verse{24}
{
	\hP{דַּבֵּר אֶל בְּנֵי יִשְׂרָאֵל לֵאמֹר בַּחֹדֶשׁ הַשְּׁבִיעִי בְּאֶחָד לַחֹדֶשׁ יִהְיֶה לָכֶם שַׁבָּתוֹן זִכְרוֹן תְּרוּעָה מִקְרָא קֹדֶשׁ׃}
}
{
	\eP{
		“Speak to the children of Israel, saying: In the seventh
		month, on the first of the month you shall have a ceasing,
		a commemoration with horn-blasting, a holy assembly.
	}
}
{
%	\footnote{
%		In the seventh month, on the first of the month. This holiday is not
%		named; but, coming nine days before Yom Kippur and involving the
%		blasting of a horn, it is clearly the holiday that is now known as the
%		Jewish New Year, Rosh Hashanah. It is not known by this name in the
%		Torah because, as this verse indicates, it came in the seventh month of
%		the calendar. In the Torah’s calendar, the first month comes in spring,
%		and the first holiday is Passover (see Exod 12:2). Footnote 23:24.
%		blasting. Sounding a shofar or trumpet.
%	}
}

\Verse{25}
{
	\hP{כּל מְלֶאכֶת עֲבֹדָה לֹא תַעֲשׂוּ וְהִקְרַבְתֶּם אִשֶּׁה לַיהֹוָה׃}
}
{
	\eP{
		You shall not do any act of work. And you shall bring
		forward an offering by fire to YHWH.”
	}
}
{
}

\Verse{26}
{
	\hP{וַיְדַבֵּר יְהֹוָה אֶל מֹשֶׁה לֵּאמֹר׃}
}
{
	\eP{And YHWH spoke to Moses, saying,}
}
{
}

\Verse{27}
{
	\hP{אַךְ בֶּעָשׂוֹר לַחֹדֶשׁ הַשְּׁבִיעִי הַזֶּה יוֹם הַכִּפֻּרִים הוּא מִקְרָא קֹדֶשׁ יִהְיֶה לָכֶם וְעִנִּיתֶם אֶת נַפְשֹׁתֵיכֶם וְהִקְרַבְתֶּם אִשֶּׁה לַיהֹוָה׃}
}
{
	\eP{
		“Just: On the tenth of this seventh month, it is the Day
		of Atonement. You shall have a holy assembly, and you
		shall degrade yourselves. And you shall bring forward an
		offering by fire to YHWH.
	}
}
{
%	\footnote{
%		Just. This term of special emphasis (Hebrew ’ak) introduces commands
%		concerning certain holidays: the Sabbath (Exod 31:13), Passover (12:16),
%		the Day of Atonement (Lev 23:27), and Sukkot (23:39). Footnote 23:27.
%		degrade yourselves. The word “degrade” is the same term that is used for
%		what the Egyptians did to the Israelites (Exod 1:11–12). As a noun or
%		adjective, the word means “poor.” As a verb, it means “to abuse,
%		afflict, oppress.” It suggests that on the Day of Atonement one is to be
%		hard on oneself and not to give care to making oneself look attractive
%		or materially successful. It is understood to include fasting.
%	}
}

\Verse{28}
{
	\hP{וְכל מְלָאכָה לֹא תַעֲשׂוּ בְּעֶצֶם הַיּוֹם הַזֶּה כִּי יוֹם כִּפֻּרִים הוּא לְכַפֵּר עֲלֵיכֶם לִפְנֵי יְהֹוָה אֱלֹהֵיכֶם׃}
}
{
	\eP{
		And you shall not do any work on this very day, because it
		is a day of atonement, to atone for you in front of YHWH,
		your {\God}.
	}
}
{
}

\Verse{29}
{
	\hP{כִּי כל הַנֶּפֶשׁ אֲשֶׁר לֹא תְעֻנֶּה בְּעֶצֶם הַיּוֹם הַזֶּה וְנִכְרְתָה מֵעַמֶּיהָ׃}
}
{
	\eP{
		Because any person who will not be degraded on this very
		day will be cut off from his people.
	}
}
{
%	\footnote{
%		his people. The Hebrew word for “person” or “soul” (nepeš) is feminine,
%		and so this is a case in which the feminine form stands for both men and
%		women—contrary to the usual practice in Hebrew (and English). A more
%		neutral English translation could be “I shall destroy that soul from
%		among its people.” But the word “soul” carries a wide range of meanings,
%		which go beyond the point of this verse. And it is worth noting that
%		there are occasions when the feminine as well as the masculine is used
%		to stand for all humans.
%	}
}

\Verse{30}
{
	\hP{וְכל הַנֶּפֶשׁ אֲשֶׁר תַּעֲשֶׂה כּל מְלָאכָה בְּעֶצֶם הַיּוֹם הַזֶּה וְהַאֲבַדְתִּי אֶת הַנֶּפֶשׁ הַהִוא מִקֶּרֶב עַמָּהּ׃}
}
{
	\eP{
		And any person who will do any work in this very day: I
		shall destroy that person from among his people.
	}
}
{
}

\Verse{31}
{
	\hP{כּל מְלָאכָה לֹא תַעֲשׂוּ חֻקַּת עוֹלָם לְדֹרֹתֵיכֶם בְּכֹל מֹשְׁבֹתֵיכֶם׃}
}
{
	\eP{
		You shall not do any work: an eternal law through your
		generations in all your homes.
	}
}
{
}

\Verse{32}
{
	\hP{שַׁבַּת שַׁבָּתוֹן הוּא לָכֶם וְעִנִּיתֶם אֶת נַפְשֹׁתֵיכֶם בְּתִשְׁעָה לַחֹדֶשׁ בָּעֶרֶב מֵעֶרֶב עַד עֶרֶב תִּשְׁבְּתוּ שַׁבַּתְּכֶם׃}
}
{
	\eP{
		It is a Sabbath, a ceasing, for you, and you shall degrade
		yourselves: on the ninth of the month in the evening, from
		evening to evening, you shall keep your Sabbath.”
	}
}
{
}

\Verse{33}
{
	\hP{וַיְדַבֵּר יְהֹוָה אֶל מֹשֶׁה לֵּאמֹר׃}
}
{
	\eP{And YHWH spoke to Moses, saying,}
}
{
}

\Verse{34}
{
	\hP{דַּבֵּר אֶל בְּנֵי יִשְׂרָאֵל לֵאמֹר בַּחֲמִשָּׁה עָשָׂר יוֹם לַחֹדֶשׁ הַשְּׁבִיעִי הַזֶּה חַג הַסֻּכּוֹת שִׁבְעַת יָמִים לַיהֹוָה׃}
}
{
	\eP{
		“Speak to the children of Israel, saying: On the fifteenth
		day of this seventh month is the Festival of Booths, seven
		days, for YHWH.
	}
}
{
%	\footnote{
%		Booths. Hebrew sukkôt.This term is used to describe the function of the
%		prket, the pavilion that covers over the ark (Exod 40:3; Ps 27:5; Lam
%		2:6; see my comment on Exod 26:31). It thus reflects simultaneously both
%		the simple, makeshift dwellings that house the Israelites in the
%		wilderness period and the beautiful, most sacred structure that houses
%		the ark.
%	}
}

\Verse{35}
{
	\hP{בַּיּוֹם הָרִאשׁוֹן מִקְרָא קֹדֶשׁ כּל מְלֶאכֶת עֲבֹדָה לֹא תַעֲשׂוּ׃}
}
{
	\eP{
		On the first day is a holy assembly. You shall not do any
		act of work.
	}
}
{
}

\Verse{36}
{
	\hP{שִׁבְעַת יָמִים תַּקְרִיבוּ אִשֶּׁה לַיהֹוָה בַּיּוֹם הַשְּׁמִינִי מִקְרָא קֹדֶשׁ יִהְיֶה לָכֶם וְהִקְרַבְתֶּם אִשֶּׁה לַיהֹוָה עֲצֶרֶת הִוא כּל מְלֶאכֶת עֲבֹדָה לֹא תַעֲשׂוּ׃}
}
{
	\eP{
		For seven days you shall bring forward an offering by fire
		to YHWH. On the eighth day you shall have a holy assembly,
		and you shall bring forward an offering by fire to YHWH.
		It is a convocation. You shall not do any act of work.
	}
}
{
%	\footnote{
%		convocation. The meaning of this term, Hebrew ‘eret, is uncertain. It
%		occurs six times in the Tanak, and five of them refer specifically to
%		this holiday, now known as Shemini Atzeret. The sixth (Jer 9:1) suggests
%		that it means a gathering of people. It is thus understood to mean a
%		special assembly.
%	}
}

\Verse{37}
{
	\hP{אֵלֶּה מוֹעֲדֵי יְהֹוָה אֲשֶׁר תִּקְרְאוּ אֹתָם מִקְרָאֵי קֹדֶשׁ לְהַקְרִיב אִשֶּׁה לַיהֹוָה עֹלָה וּמִנְחָה זֶבַח וּנְסָכִים דְּבַר יוֹם בְּיוֹמוֹ׃}
}
{
	\eP{
		“These are YHWH’s appointed times, which you shall call
		holy assemblies, to bring forward an offering by fire to
		YHWH: burnt offering and grain offering, sacrifice and
		libations, each day’s thing on its day,
	}
}
{
}

\Verse{38}
{
	\hP{מִלְּבַד שַׁבְּתֹת יְהֹוָה וּמִלְּבַד מַתְּנוֹתֵיכֶם וּמִלְּבַד כּל נִדְרֵיכֶם וּמִלְּבַד כּל נִדְבֹתֵיכֶם אֲשֶׁר תִּתְּנוּ לַיהֹוָה׃}
}
{
	\eP{
		aside from YHWH’s Sabbaths and aside from your gifts and
		aside from all of your vows and aside from all of your
		contributions that you will give to YHWH.
	}
}
{
}

\Verse{39}
{
	\hP{אַךְ בַּחֲמִשָּׁה עָשָׂר יוֹם לַחֹדֶשׁ הַשְּׁבִיעִי בְּאסְפְּכֶם אֶת תְּבוּאַת הָאָרֶץ תָּחֹגּוּ אֶת חַג יְהֹוָה שִׁבְעַת יָמִים בַּיּוֹם הָרִאשׁוֹן שַׁבָּתוֹן וּבַיּוֹם הַשְּׁמִינִי שַׁבָּתוֹן׃}
}
{
	\eP{
		“Just: On the fifteenth day of the seventh month when you
		gather the land’s produce, you shall celebrate YHWH’s
		holiday seven days. On the first day is a ceasing, and on
		the eighth day is a ceasing.
	}
}
{
}

\Verse{40}
{
	\hP{וּלְקַחְתֶּם לָכֶם בַּיּוֹם הָרִאשׁוֹן פְּרִי עֵץ הָדָר כַּפֹּת תְּמָרִים וַעֲנַף עֵץ עָבֹת וְעַרְבֵי נָחַל וּשְׂמַחְתֶּם לִפְנֵי יְהֹוָה אֱלֹהֵיכֶם שִׁבְעַת יָמִים׃}
}
{
	\eP{
		And on the first day you shall take fruit of appealing
		trees, branches of palms, boughs of thick trees, and
		willows of a wadi, and you shall be happy in front of
		YHWH, your {\God}, seven days.
	}
}
{
}

\Verse{41}
{
	\hP{וְחַגֹּתֶם אֹתוֹ חַג לַיהֹוָה שִׁבְעַת יָמִים בַּשָּׁנָה חֻקַּת עוֹלָם לְדֹרֹתֵיכֶם בַּחֹדֶשׁ הַשְּׁבִיעִי תָּחֹגּוּ אֹתוֹ׃}
}
{
	\eP{
		And you shall celebrate it, a holiday to YHWH, seven days
		in the year, an eternal law through your generations; you
		shall celebrate it in the seventh month.
	}
}
{
}

\Verse{42}
{
	\hP{בַּסֻּכֹּת תֵּשְׁבוּ שִׁבְעַת יָמִים כּל הָאֶזְרָח בְּיִשְׂרָאֵל יֵשְׁבוּ בַּסֻּכֹּת׃}
}
{
	\eP{
		You shall live in booths seven days. Every citizen in
		Israel shall live in booths,
	}
}
{
%	\footnote{
%		booths. Hebrew sukkôt. See comment on 23:34.
%	}
}

\Verse{43}
{
	\hP{לְמַעַן יֵדְעוּ דֹרֹתֵיכֶם כִּי בַסֻּכּוֹת הוֹשַׁבְתִּי אֶת בְּנֵי יִשְׂרָאֵל בְּהוֹצִיאִי אוֹתָם מֵאֶרֶץ מִצְרָיִם אֲנִי יְהֹוָה אֱלֹהֵיכֶם׃}
}
{
	\eP{
		so your generations will know that I had the children of
		Israel live in booths when I brought them out from the
		land of Egypt. I am YHWH, your {\God}.”
	}
}
{
}

\Verse{44}
{
	\hP{וַיְדַבֵּר מֹשֶׁה אֶת מֹעֲדֵי יְהֹוָה אֶל בְּנֵי יִשְׂרָאֵל׃}
}
{
	\eP{
		And Moses spoke YHWH’s appointed times to the children of
		Israel.
	}
}
{
}
